\documentclass[a4paper,10pt]{article}
\usepackage[utf8]{inputenc}
\usepackage[T1]{fontenc}
\usepackage[english]{babel}
\usepackage[a4paper, landscape, top=1cm, bottom=1cm, left=1cm, right=1cm]{geometry}
\usepackage{tabularx}
\usepackage{array}
\usepackage{xfp}
\usepackage[table]{xcolor}
\usepackage{helvet}
\renewcommand{\familydefault}{\sfdefault} % sans-serif reads better at small sizes in dense tables
\usepackage{titlesec}
\usepackage{enumitem}

% =========================================================
% VISUAL IDENTITY
% =========================================================
\definecolor{accent}{HTML}{34495E}   % neutral heading/rule color (Testing Protocol, Appendix)
\definecolor{ssaccent}{HTML}{B8860B} % Phase 1 heading tint (Sweet Spot gold)
\definecolor{z4accent}{HTML}{D2691E} % Phase 2 heading tint (Threshold orange)
\definecolor{z5accent}{HTML}{C0334D} % Phase 3 heading tint (VO2 red)

\titleformat{\section}{\large\bfseries}{}{0em}{}[{\color{accent}\titlerule[0.8pt]}]
\titlespacing*{\section}{0pt}{2pt}{6pt}
\titleformat{\subsection}{\normalsize\bfseries\color{accent}}{}{0em}{}
\titlespacing*{\subsection}{0pt}{0pt}{4pt}

% Callout box for phase rationale / protocol notes — replaces plain italic paragraphs
\newcommand{\NoteBox}[1]{%
  \noindent\colorbox{accent!6}{\begin{minipage}{\dimexpr\textwidth-2\fboxsep\relax}
  \footnotesize #1
  \end{minipage}}%
}

% =========================================================
% USER SETTINGS — update CP/W' and Threshold Pace here after each retest
% =========================================================
\def\currentCP{260}     % cycling Critical Power in watts (see Testing Protocol)
\def\currentWprime{20}  % anaerobic work capacity W' in kJ, from the same CP test
\def\currentCS{280}     % running Threshold Pace in seconds per km
                        %   240 = 4:00/km, 255 = 4:15/km, 270 = 4:30/km
                        % Estimate via 3-min all-out test, or two TTs
                        % (e.g. 3 min + 9 min): TP = (D2 - D1) / (T2 - T1).

% =========================================================
% ZONE COLORS
% =========================================================
\definecolor{cz1}{HTML}{A8D8EA}      % Z1  Recovery
\definecolor{cz2}{HTML}{5AB49E}      % Z2  Endurance
\definecolor{cz3}{HTML}{F4C36A}      % Z3  Tempo (used by run cells)
\definecolor{css}{HTML}{FCD663}      % SS  Sweet Spot
\definecolor{cz4}{HTML}{FE9C5B}      % Z4  Threshold
\definecolor{cz5}{HTML}{E65476}      % Z5  VO2 Max
\definecolor{cz6}{HTML}{8E44AD}      % Z6  Anaerobic (reference table only)
\definecolor{recweek}{HTML}{E8E8E8}  % Recovery week rows

% =========================================================
% DYNAMIC HELPERS — CP (cycling) and Threshold Pace (running)
% =========================================================

% Watts at given % CP: \W{90} -> "230"
\newcommand{\W}[1]{\fpeval{round(\currentCP*#1/100,0)}}

% Single % CP with watts: \sgl{106} -> "106% (270W)"
\newcommand{\sgl}[1]{#1\%~{\scriptsize(\W{#1}W)}}

% Range % CP with watts: \rng{90}{94} -> "90-94% (230-240W)"
\newcommand{\rng}[2]{#1--#2\%~{\scriptsize(\W{#1}--\W{#2}W)}}

% Run pace at % of Threshold Pace speed: \rpace{82} -> "4:53/km"
% Total seconds computed once into \rp@s to avoid triple re-evaluation.
\newcommand{\rpace}[1]{%
  \edef\rp@s{\fpeval{round(\currentCS/(#1/100),0)}}%
  \fpeval{floor(\rp@s/60)}:%
  \ifnum\fpeval{\rp@s-60*floor(\rp@s/60)}<10 0\fi%
  \fpeval{\rp@s-60*floor(\rp@s/60)}/km%
}

% =========================================================
% CELL BACKGROUND SHORTHANDS
% =========================================================
\newcommand{\hlZone}{\cellcolor{cz1!35}}
\newcommand{\hlZtwo}{\cellcolor{cz2!35}}
\newcommand{\hlZthree}{\cellcolor{cz3!35}}
\newcommand{\hlSS}{\cellcolor{css!35}}
\newcommand{\hlZfour}{\cellcolor{cz4!35}}
\newcommand{\hlZfive}{\cellcolor{cz5!35}}

% Auto-select run cell color from % of Threshold Pace speed.
% Boundaries match Garmin/Lactate zone image:
%   <78   Z1  Recovery        -> cz1   (0-77.5%)
%   78-87 Z2  Endurance       -> cz2   (78.5-87.7%)
%   88-94 Z3  Tempo           -> cz3   (88.7-94.3%)
%   95-100 Z4 Threshold       -> cz4   (95.3-100%)
%   >=101 Z5+ VO2/Anaerobic   -> cz5   (101%+)
\newcommand{\runcolor}[1]{%
  \ifnum#1<78\hlZone\else
  \ifnum#1<88\hlZtwo\else
  \ifnum#1<95\hlZthree\else
  \ifnum#1<101\hlZfour\else
  \hlZfive\fi\fi\fi\fi
}

% =========================================================
% PER-DAY CELL MACROS — one training per column
% =========================================================

% WO1 — recovery spin (always identical)
\newcommand{\DayMon}{%
  \hlZone\footnotesize\textbf{Z1 Spin}\par
  {\scriptsize 90'\,@\,\rng{45}{50}}%
}

% Bike workout cell: \DayBike{\hlSS}{title}{set spec}{rec / note}
\newcommand{\DayBike}[4]{%
  #1\footnotesize\textbf{#2}\par
  {\scriptsize #3\par\textit{#4}}%
}

% Run cell: \DayRun{label}{km}{lo%TP}{hi%TP}
% Background color driven by lo% (ranges stay within one zone).
\newcommand{\DayRun}[4]{%
  \runcolor{#3}\footnotesize\textbf{#1 Run}\par
  {\scriptsize #2\,km @ #3--#4\%~TP\par(\rpace{#4}--\rpace{#3})}%
}

% Tempo interval run cell: \DayRunInt{sets}{duration}{lo%TP}{hi%TP}{rec}
% Color uses hi% so threshold-range work renders as Z4, not Z3.
\newcommand{\DayRunInt}[5]{%
  \runcolor{#4}\footnotesize\textbf{Tempo Intervals}\par
  {\scriptsize #1$\times$#2' @ #3--#4\%~TP\par(\rpace{#4}--\rpace{#3})\par\textit{#5 jog rec}}%
}

% Weekend ride cell: \DayWE{day label}{hours}{lo%}{hi%}
\newcommand{\DayWE}[4]{%
  \hlZtwo\footnotesize\textbf{#1}\par
  {\scriptsize #2\,h @ #3--#4\%\par(\W{#3}--\W{#4}W)}%
}

% Recovery-week row spanning all 7 day columns
\newcommand{\RecoveryRow}[2]{%
  \textbf{#1} &
  \multicolumn{7}{c|}{\cellcolor{recweek}\rule{0pt}{1.4cm}%
    \textbf{RECOVERY WEEK --- #2}}\\
  \hline
}

% =========================================================
% TABLE GEOMETRY
% =========================================================
\newcolumntype{Y}{>{\raggedright\arraybackslash}X}
\renewcommand{\arraystretch}{1.4}

% =========================================================
% DOCUMENT
% =========================================================
\title{\textbf{\Huge Hybrid Aerobic Development Plan}\\[0.3em]
       \large Base $\to$ Threshold $\to$ VO\textsubscript{2}Max \quad (12 Weeks)}
\author{CP: \currentCP{}\,W ($W'$ \currentWprime{}\,kJ) \quad\textbullet\quad Threshold Pace: \rpace{100} \\[0.2em]
        \textit{\small Update \texttt{\textbackslash currentCP}, \texttt{\textbackslash currentWprime} and
                       \texttt{\textbackslash currentCS} --- all watts and paces recompute. See Testing Protocol.}}
\date{\today}

\begin{document}
\maketitle


% =========================================================
% PHASE 1 — BASE
% =========================================================
\section*{\color{ssaccent!70!black}Phase 1 --- Base: Sweet Spot \& Volume (Weeks 1--4)}

\NoteBox{Sweet Spot, not pure Zone 2, is the primary quality tool this phase. Zone 2 isn't
intrinsically superior for raising thresholds or VO\textsubscript{2}max --- at matched work, zone
3--5 training drives larger threshold/VO\textsubscript{2}max gains per unit of effort (Inglis et al.
2024), and low-intensity volume is a complement to 2--3 weekly quality sessions, not a mandatory
prerequisite before intervals (Matomäki 2025). Sweet spot is preferred over straight threshold here
because it produces near-identical adaptations at meaningfully lower fatigue cost, letting SS Quality
(WO2) and SS Long (WO4) coexist in the same week without compounding into an unrecoverable load.
Progression is dosed as time-in-zone, not power: SS Quality adds a set/duration each week
(2$\times$20$\to$2$\times$25$\to$3$\times$20) while SS Long extends continuously
(40$\to$45$\to$50\,min) --- staying under the $\sim$2--2.5\,h per-session ceiling and the 2$\times$/week
frequency guardrail past which sweet-spot fatigue starts to outpace its signal. Endurance runs and
weekend Z2 rides progress the same way, duration before intensity, since extending duration is the
correct lever until race-relevant duration is covered. \textbf{Before W1:} run the Critical Power /
$W'$ test and the threshold-pace test in the Testing Protocol section --- don't reuse an old or
estimated number; an overestimated CP turns ``SS Quality'' (90\%) into threshold from day one.}

\medskip

\noindent
\begin{tabularx}{\textwidth}{|>{\centering\arraybackslash}p{0.5cm}|*{7}{Y|}}
\hline
\rowcolor{black!10}
\textbf{W} & \textbf{WO1} & \textbf{WO2} & \textbf{WO3} & \textbf{WO4} & \textbf{WO5} & \textbf{WO6} & \textbf{WO7} \\
\hline

% W1
\textbf{1} &
\DayMon &
\DayBike{\hlSS}{SS Quality}{2x20' @ \sgl{90}}{4' rec, 85--95 rpm, baseline} &
\DayRunInt{3}{12}{89}{94}{2'} &
\DayBike{\hlSS}{SS Long}{1x40' @ \rng{88}{90}}{continuous, TTE anchor} &
\DayRun{Endurance}{10}{79}{87} &
\DayWE{WO6 Long Z2}{3}{65}{72} &
\DayWE{WO7 Med Z2}{2}{62}{70}
\\ \hline

% W2
\textbf{2} &
\DayMon &
\DayBike{\hlSS}{SS Quality}{2x25' @ \sgl{90}}{4' rec, +5' per interval} &
\DayRunInt{3}{14}{89}{94}{2'} &
\DayBike{\hlSS}{SS Long}{1x45' @ \rng{88}{90}}{continuous, +5' vs W1} &
\DayRun{Endurance}{12}{79}{87} &
\DayWE{WO6 Long Z2}{3.5}{65}{72} &
\DayWE{WO7 Med Z2}{2}{62}{70}
\\ \hline

% W3
\textbf{3} &
\DayMon &
\DayBike{\hlSS}{SS Quality}{3x20' @ \sgl{90}}{4' rec, +1 set vs W1} &
\DayRunInt{2}{25}{89}{94}{2'} &
\DayBike{\hlSS}{SS Long}{1x50' @ \rng{88}{90}}{continuous, +5' vs W2} &
\DayRun{Endurance}{14}{79}{87} &
\DayWE{WO6 Long Z2}{4}{63}{72} &
\DayWE{WO7 Med Z2}{2.5}{62}{70}
\\ \hline

% W4 RECOVERY
\RecoveryRow{4}{Self-Managed}

\end{tabularx}

\newpage

% =========================================================
% PHASE 2 — THRESHOLD
% =========================================================
\section*{\color{z4accent!70!black}Phase 2 --- Threshold: CP Ceiling Development (Weeks 5--8)}

\NoteBox{Threshold succeeds sweet spot once SS Long's TTE growth signals the natural next step: W5's
Threshold Entry (3$\times$12' @ 95--100\%) starts conservative and extends toward TTE
(3$\times$20' by W7), following the standard threshold progressive-overload rule --- add
time-in-zone per session, not more sets. The second weekly quality day (WO4) uses over/unders
instead of a second steady-state threshold session: O/Us carry no unique physiological edge over
steady-state work, but substituting one threshold day for O/Us adds race-specific accelerate-and-recover
practice and keeps the week varied \emph{without} adding load on top of the existing threshold day.
Overs sit at 108\%\ CP ($\sim$+20\,W) rather than the commonly-misapplied 120\%\ --- that range
crosses into VO\textsubscript{2}max territory and defeats the point of a threshold session; unders
sit at sweet spot (90\%\ CP), matching the standard 1:4 bread-and-butter format (1\,min over /
4\,min under) and its documented progression (3$\times$15$\to$3$\times$20$\to$3$\times$25\,min).
W8's recovery and CP/TP retest closes the phase and is the decision gate: if threshold TTE and CP
are still moving, more threshold work is still the right tool; if they've stalled, that's the signal
to move to VO\textsubscript{2} work (Phase 3) rather than add more threshold volume.}

\medskip

\noindent
\begin{tabularx}{\textwidth}{|>{\centering\arraybackslash}p{0.5cm}|*{7}{Y|}}
\hline
\rowcolor{black!10}
\textbf{W} & \textbf{WO1} & \textbf{WO2} & \textbf{WO3} & \textbf{WO4} & \textbf{WO5} & \textbf{WO6} & \textbf{WO7} \\
\hline

% W5
\textbf{5} &
\DayMon &
\DayBike{\hlZfour}{Threshold Entry}{3x12' @ \rng{95}{100}}{3' rec, start conservative} &
\DayRunInt{3}{12}{89}{94}{2'} &
\DayBike{\hlZfour}{O/U Short}{3x15' (1' \sgl{108} / 4' \sgl{90})}{5' rec} &
\DayRun{Endurance}{14}{79}{87} &
\DayWE{WO6 Long Z2}{3.5}{65}{72} &
\DayWE{WO7 Med Z2}{2}{62}{70}
\\ \hline

% W6
\textbf{6} &
\DayMon &
\DayBike{\hlZfour}{Threshold Build}{3x15' @ \rng{95}{100}}{3' rec, even power} &
\DayRunInt{3}{14}{89}{94}{2'} &
\DayBike{\hlZfour}{O/U Medium}{3x20' (1' \sgl{108} / 4' \sgl{90})}{5' rec} &
\DayRun{Endurance}{16}{79}{87} &
\DayWE{WO6 Long Z2}{3.5}{65}{72} &
\DayWE{WO7 Med Z2}{2}{62}{70}
\\ \hline

% W7
\textbf{7} &
\DayMon &
\DayBike{\hlZfour}{TTE Extension}{3x20' @ \rng{95}{100}}{5' rec, continuous preferred if legs} &
\DayRunInt{2}{25}{89}{94}{2'} &
\DayBike{\hlZfour}{O/U Long}{3x25' (1' \sgl{108} / 4' \sgl{90})}{5' rec, peak threshold} &
\DayRun{Endurance}{18}{79}{87} &
\DayWE{WO6 Long Z2}{4}{63}{72} &
\DayWE{WO7 Med Z2}{2.5}{62}{70}
\\ \hline

% W8 RECOVERY + RETEST
\RecoveryRow{8}{Self-Managed + CP/TP Retest}

\end{tabularx}

\newpage

% =========================================================
% PHASE 3 — VO2 MAX (Classical Distributed Periodization, Empirical Cycling protocol)
% =========================================================
\section*{\color{z5accent!70!black}Phase 3 --- VO\textsubscript{2}Max: Progressive Development (Weeks 9--12)}

\NoteBox{Two VO\textsubscript{2} sessions/week, not a concentrated block. A 5-session/week
HIT block (Rønnestad et al. 2014) was considered and dropped: it compounds two variables --- session
frequency and per-rep intensity --- that no cited study tested together, and it relies on a volume
cutback and elite-level recovery capacity this plan doesn't assume. Progressive overload instead runs
\textbf{within} the 2x/week format across the block: reps climb 4$\to$5$\to$6 and recovery trims
3$\to$2.5$\to$2\,min week to week, while total work per session stays inside the 15--25\,min ceiling
--- volume is not what's overloaded here, effort quality is.
Per-interval execution follows the Empirical Cycling protocol: \textbf{no fixed power target} --- start
each rep 20--30\% harder than you think you can hold, let power fade, keep pushing on perceived effort
until breathing can no longer stay maxed. \textbf{Cadence 100--110\,rpm} throughout --- cadence, not
power, is the primary lever for cardiac preload (stroke volume), which is the actual adaptation target.
Recovery between reps: start shorter than instinct, add time only if the next rep can't reach max again.
W9--11 carry full LIT volume alongside the 2 VO\textsubscript{2} days (easy run replaces tempo-interval
running for this phase, protecting recovery around the two hard bike days). W12 is recovery; retest at
the END of the week, not the start --- VO\textsubscript{2}max gains surface with a delay.}

\medskip

\noindent
\begin{tabularx}{\textwidth}{|>{\centering\arraybackslash}p{0.5cm}|*{7}{Y|}}
\hline
\rowcolor{black!10}
\textbf{W} & \textbf{WO1} & \textbf{WO2} & \textbf{WO3} & \textbf{WO4} & \textbf{WO5} & \textbf{WO6} & \textbf{WO7} \\
\hline

% W9 — VO2 entry (4 reps, longer recovery)
\textbf{9} &
\DayMon &
\DayBike{\hlZfive}{VO2 A}{4x4' max mean / 100--110 rpm}{3' rec, breathing maxed last 60--90s} &
\DayRun{Easy}{12}{65}{77} &
\DayWE{WO4 Z2}{1.5}{60}{68} &
\DayBike{\hlZfive}{VO2 B}{4x5' max mean / 100--110 rpm}{3' rec, let power fade, hold effort} &
\DayWE{WO6 Long Z2}{3.5}{65}{72} &
\DayWE{WO7 Med Z2}{2}{62}{70}
\\ \hline

% W10 — VO2 build (5 reps, recovery trimmed)
\textbf{10} &
\DayMon &
\DayBike{\hlZfive}{VO2 A}{5x4' max mean / 100--110 rpm}{2.5' rec, breathing maxed last 60--90s} &
\DayRun{Easy}{14}{65}{77} &
\DayWE{WO4 Z2}{1.5}{60}{68} &
\DayBike{\hlZfive}{VO2 B}{4x5' max mean / 100--110 rpm}{2.5' rec, let power fade, hold effort} &
\DayWE{WO6 Long Z2}{4}{63}{72} &
\DayWE{WO7 Med Z2}{2.5}{62}{70}
\\ \hline

% W11 — VO2 peak (6 reps, shortest recovery, converges to target volume)
\textbf{11} &
\DayMon &
\DayBike{\hlZfive}{VO2 A}{6x4' max mean / 100--110 rpm}{2' rec, breathing maxed last 60--90s} &
\DayRun{Easy}{16}{65}{77} &
\DayWE{WO4 Z2}{1.5}{60}{68} &
\DayBike{\hlZfive}{VO2 B}{5x5' max mean / 100--110 rpm}{2' rec, let power fade, hold effort} &
\DayWE{WO6 Long Z2}{4}{63}{72} &
\DayWE{WO7 Med Z2}{2.5}{62}{70}
\\ \hline

% W12 RECOVERY — delayed-effect window
\RecoveryRow{12}{Adaptation surfaces; CP/TP retest at END of week}

\end{tabularx}

\newpage
% =========================================================
% APPENDIX: ZONE REFERENCES (CP + Threshold Pace side by side)
% =========================================================
\section*{Appendix --- Zone Reference}

\begin{minipage}[t]{0.34\textwidth}
\subsection*{Cycling Zones at CP = \currentCP{}\,W}
\footnotesize
\begin{tabular}{@{}lll@{}}
\textbf{Zone} & \textbf{\% CP} & \textbf{Watts} \\
\hline
\rowcolor{cz1!35}     Z1 Recovery  & $<$56\%      & $<$\W{56}W \\
\rowcolor{cz2!35}     Z2 Endurance & 56--75\%     & \W{56}--\W{75}W \\
\rowcolor{cz3!35}     Z3 Tempo     & 76--87\%     & \W{76}--\W{87}W \\
\rowcolor{css!35}     Sweet Spot   & 88--94\%     & \W{88}--\W{94}W \\
\rowcolor{cz4!35}     Z4 Threshold & 95--105\%    & \W{95}--\W{105}W \\
\rowcolor{cz5!35}     Z5 VO2 Max   & 106--120\%   & \W{106}--\W{120}W \\
\rowcolor{cz6!35}     Z6 Anaerobic & $>$121\%     & $>$\W{121}W \\
\end{tabular}
\end{minipage}
\hspace{0.06\textwidth}%
\begin{minipage}[t]{0.40\textwidth}
\subsection*{Running Zones at Threshold Pace = \rpace{100}}
\footnotesize
\begin{tabular}{@{}lll@{}}
\textbf{Zone} & \textbf{\% TP} & \textbf{Pace} \\
\hline
\rowcolor{cz1!35} Z1 Recovery       & $<$77.5\%      & slower than \rpace{78} \\
\rowcolor{cz2!35} Z2 Endurance      & 78.5--87.7\%   & \rpace{88}--\rpace{79} \\
\rowcolor{cz3!35} Z3 Tempo          & 88.7--94.3\%   & \rpace{94}--\rpace{89} \\
\rowcolor{cz4!35} Z4 Threshold      & 95.3--100\%    & \rpace{100}--\rpace{95} \\
\rowcolor{cz5!25} Z5a VO\textsubscript{2}Max & 101--103.4\%  & \rpace{103}--\rpace{101} \\
\rowcolor{cz5!45} Z5b Anaerobic     & 104.4--111.5\% & \rpace{112}--\rpace{104} \\
\rowcolor{cz5!65} Z5c Neuromuscular & $>$112.5\%     & faster than \rpace{113} \\
\end{tabular}
\end{minipage}

\bigskip

% =========================================================
% TESTING PROTOCOL — Critical Power & W' (replaces the FTP-95%-of-20min shortcut)
% =========================================================
\section*{Testing Protocol --- Critical Power (CP) \& $W'$}

\NoteBox{This plan anchors all zones to \textbf{Critical Power (CP)}, not the classic
``95\%\ of 20-min'' FTP shortcut. Head-to-head studies (Morgan 2019, Karsten 2021) show FTP and CP agree
on average across a group but diverge by up to 15\%\ for a given individual --- enough to turn an intended
Sweet Spot session into unintentional threshold work. CP is the field-testable proxy for the Maximum
Metabolic Steady State (the boundary between sustainable and rapidly-fatiguing effort), derived from the
power--duration relationship across multiple maximal efforts rather than a single 20-min time trial.}

\medskip

\noindent\textbf{\footnotesize Protocol --- 3 maximal efforts, all-out:}
\begin{itemize}[nosep, leftmargin=1.4em, itemsep=2pt, label=\textbf{\textcolor{accent}{--}}]
\footnotesize
  \item \textbf{Efforts:} 3\,min, 8\,min, 12--13\,min --- all above CP (severe-intensity domain). Not
        shorter (risks hitting VO\textsubscript{2}max/exhaustion before the effort is even representative)
        and not longer (glycogen, pacing, and motivation start dominating the result).
  \item \textbf{Recovery between efforts:} $\geq$30\,min if done in one session, or spread across separate
        days ($\geq$24\,h) to remove fatigue carry-over between trials --- prefer separate days if your
        schedule allows it.
  \item \textbf{Consistency:} same terrain across all three efforts (flat or a steady climb --- whichever
        you train/race on most), same bike position, cadence held within $\pm$5\,rpm.
  \item \textbf{Pacing:} controlled fast start ($\sim$10\%\ above goal power for the first 30\,s, not
        all-out), even effort through the middle, a final 30\,s kick if you have it left. Pick a route free
        of lights, traffic, and gradient changes.
  \item \textbf{Fitting CP and $W'$:} plot power vs.\ duration (or total work vs.\ duration) and fit all
        three standard models --- 2-parameter hyperbolic, linear power-vs-1/time, linear work-vs-time. Take
        the \textbf{best individual fit}: whichever model gives the lowest combined error (SEE) for CP and
        $W'$ on \emph{your} data, not a fixed default model.
  \item \textbf{Quality check:} if CP's SEE exceeds 5\%\ or $W'$'s SEE exceeds 10\%\ across all three
        models, re-run whichever trial looks like it was paced poorly rather than trusting the fit.
\end{itemize}

\medskip

\NoteBox{Update \texttt{\textbackslash currentCP} and \texttt{\textbackslash currentWprime}
with the result. Retest at the checkpoints already built into this plan: end of W8 (recovery week) and end
of W12 (after the VO\textsubscript{2}max adaptation window closes, not before). For running Threshold Pace,
estimate the same way with two to three all-out time trials (e.g.\ 3\,min + 9--13\,min) on a flat, consistent
route: $TP = (D_2 - D_1)/(T_2 - T_1)$.}

\end{document}